% =============================================================================
% Problem Template
% =============================================================================
% Copy this file into:
%   latex/sections/<domain>/problems/lc-XXX-problem-name.tex
%
% Fill in the metadata below and remove any sections not required.
% =============================================================================

% STATUS: draft | reviewed | final
% PROMOTE: none | exemplars | quant-prep
% TAGS: arrays, hashing, dp, graphs, etc.
% DIFFICULTY: easy | medium | hard
% SOURCE: LeetCode XXX | Custom | Other

\ProblemTitle{Add Two Numbers}
\ProblemMeta{Medium}{Linked Lists, Math, Carry Propagation}{LC-002}
\label{prob:lc-002-add-two-numbers}

\ProblemSummary
Given two non-empty linked lists representing two non-negative integers, where
each node stores a single digit and the digits are arranged in reverse order
(least significant digit first), compute the sum of the two integers.
Return the result as a new linked list in the same reverse-order format.

\KeyIdea
Because the digits are stored in reverse order, the lists can be traversed
from least significant to most significant digit, making the problem a direct
simulation of elementary addition. At each step, the current digits and any
incoming \textbf{carry} determine the next output digit and updated carry.
A \textbf{dummy head} node anchors the result list, while a moving
\textbf{cursor} appends newly computed digits without losing access to the
list's beginning. Traversal continues until both input lists reach
\textbf{null termination} and no carry remains, maintaining the \textbf{invariant}
that all previously processed digits in the result list are correct.

% -----------------------------------------------------------------------------
% Optional: Author Thinking (excluded from exemplar builds)
% -----------------------------------------------------------------------------
\begin{Thinking}
A first brute-force idea is to traverse each list, reconstruct the two integers,
add them, and then convert the sum back into a linked list. However, this
approach is awkward in practice: the numbers may exceed native integer ranges,
and it performs extra passes and conversions that are unnecessary.

A better approach is to simulate column-wise addition directly on the linked
lists. Since the digits are stored least-significant-first, we can traverse
both lists once, propagate a carry, and build the result list incrementally.
\end{Thinking}

% -----------------------------------------------------------------------------
% Algorithm
% -----------------------------------------------------------------------------
\Algorithm
\begin{CodeBlock}[style=pseudo,caption={Pseudocode}]{16}
    dummy = new Node(0)
    cur = dummy
    carry = 0

    while l1 != null OR l2 != null OR carry != 0:
        x = (l1 != null) ? l1.val : 0
        y = (l2 != null) ? l2.val : 0

        total = x + y + carry
        digit = total mod 10
        carry = total div 10

        cur.next = new Node(digit)
        cur = cur.next

        if l1 != null: l1 = l1.next
        if l2 != null: l2 = l2.next

    return dummy.next
\end{CodeBlock}


% -----------------------------------------------------------------------------
% Correctness
% -----------------------------------------------------------------------------
\Correctness
At the start of each iteration of the loop, the result list contains the correct
sum of all digit positions processed so far, and the \textbf{carry} variable
holds the correct carry-out into the next higher digit position.
During the iteration, the algorithm adds the current digits (or zero if a list
has reached \textbf{null termination}) together with the carry, computes the
correct output digit, and updates the carry accordingly.
Because each iteration advances at least one input pointer or consumes a
remaining carry, the loop eventually terminates.
When the loop ends, all digit positions and any final carry have been processed,
so the resulting linked list represents the correct sum.


% -----------------------------------------------------------------------------
% Complexity
% -----------------------------------------------------------------------------
\Complexity
Let $n$ and $m$ be the lengths of the two input linked lists.
Each iteration processes at most one node from each list, and the loop runs for
at most $\max(n, m) + 1$ iterations, yielding a time complexity of
$O(\max(n, m))$.

The algorithm uses $O(1)$ auxiliary space, as it maintains only a constant
number of pointers and a carry variable.
The output linked list requires $O(\max(n, m))$ space to store the result.


% -----------------------------------------------------------------------------
% Implementation
% -----------------------------------------------------------------------------
\bigskip
\noindent\textbf{C++ Implementation}

\lstinputlisting[
  style=cpp,
  caption={C++ Solution: LC-002 Add Two Numbers}
]{../problems/06-linked-structures/medium/lc-002-add-two-numbers/solution.cpp}

\bigskip
\noindent\textbf{Python Implementation}

\lstinputlisting[
  language=Python,
  backgroundcolor=\color{codebg},
  basicstyle=\ttfamily\small,
  frame=single,
  breaklines=true,
  caption={Python Solution: LC-002 Add Two Numbers
]{../problems/06-linked-structures/medium/lc-002-add-two-numbers/solution.py}

\ProblemDivider
